\documentclass[11pt,a4paper]{article}

\usepackage[utf8]{inputenc}
\usepackage[T1]{fontenc}
\renewcommand{\contentsname}{Contents}
\usepackage{lmodern}
\usepackage[a4paper,margin=2.1cm,headheight=15pt]{geometry}
\usepackage{microtype}
\usepackage{graphicx}
\usepackage{booktabs}
\usepackage{tabularx}
\usepackage{array}
\usepackage{enumitem}
\usepackage{xcolor}
\usepackage{fancyhdr}
\usepackage{lastpage}
\usepackage{tcolorbox}
\usepackage{amssymb}
\usepackage{tikz}
\usetikzlibrary{arrows.meta,positioning,shapes.geometric}
\usepackage{hyperref}

% -----------------------------------------------------------------
% SENG 456 -- Individual Project Submission Form (completed)
% -----------------------------------------------------------------

\definecolor{courseNavy}{HTML}{17365D}
\definecolor{courseBlue}{HTML}{2F5597}
\definecolor{courseTeal}{HTML}{2F6F6D}
\definecolor{courseOrange}{HTML}{A65E20}
\definecolor{courseRed}{HTML}{9C2F2F}
\definecolor{coursePaper}{HTML}{F7F4EC}
\definecolor{courseGray}{HTML}{5B6573}

\hypersetup{
  colorlinks=true,
  linkcolor=courseBlue,
  urlcolor=courseBlue,
  pdfauthor={Mohammed Mustafa Kareem},
  pdftitle={SENG 456 Individual Project Submission Form -- PURPLE-LOOP},
  pdfsubject={Agent Orchestration and Multimodal Systems Term Project}
}

\setlength{\parindent}{0pt}
\setlength{\parskip}{5pt}
\setlist[itemize]{leftmargin=1.6em,itemsep=2pt,topsep=3pt}
\setlist[enumerate]{leftmargin=1.8em,itemsep=3pt,topsep=3pt}
\renewcommand{\arraystretch}{1.25}
\newcolumntype{Y}{>{\raggedright\arraybackslash}X}
\newcolumntype{C}[1]{>{\centering\arraybackslash}p{#1}}
\newcolumntype{L}[1]{>{\raggedright\arraybackslash}p{#1}}

\pagestyle{fancy}
\fancyhf{}
\fancyhead[L]{\small SENG 456 -- Individual Project Submission Form}
\fancyhead[R]{\small \textcolor{courseGray}{2026 Summer School}}
\fancyfoot[L]{\small OSTIM Technical University -- Software Engineering}
\fancyfoot[R]{\small Page \thepage/\pageref*{LastPage}}
\renewcommand{\headrulewidth}{0.4pt}
\renewcommand{\footrulewidth}{0.4pt}

\newtcolorbox{projectbox}[2][]{
  colback=coursePaper,
  colframe=#2,
  boxrule=0.8pt,
  arc=2mm,
  left=3mm,right=3mm,top=2mm,bottom=2mm,
  before skip=6pt,after skip=6pt,
  #1
}

\newtcolorbox{answerbox}{
  colback=white,
  colframe=courseGray!45,
  boxrule=0.5pt,
  arc=1.2mm,
  left=2.5mm,right=2.5mm,top=1.5mm,bottom=1.5mm,
  before skip=3pt,after skip=8pt
}

\newcommand{\sectionintro}[1]{\textcolor{courseGray}{\small #1}\par\vspace{2pt}}
\newcommand{\qa}[2]{\textbf{#1}\par\vspace{-2pt}\begin{answerbox}#2\end{answerbox}}
\newcommand{\yes}{$\checkmark$}
\newcommand{\pl}{\texttt{purple\_loop.py}}

\begin{document}
\pagenumbering{roman}

% ================================================================
% COVER
% ================================================================
\begin{titlepage}
  \thispagestyle{empty}
  \begin{center}
    \vspace*{0.2cm}
    \IfFileExists{logos/LUT-LOGO-PNG.png}{%
      \includegraphics[width=0.22\textwidth]{logos/LUT-LOGO-PNG.png}\par
    }{}
    \vspace{0.25cm}
    {\large\bfseries OSTIM TECHNICAL UNIVERSITY\par}
    {\normalsize Department of Software Engineering\par}
    \vspace{0.35cm}
    {\color{courseNavy}\rule{\linewidth}{1.4pt}\par}
    \vspace{0.15cm}
    {\Large\bfseries SENG 456\par}
    {\large\bfseries AGENT ORCHESTRATION AND MULTIMODAL SYSTEMS\par}
    \vspace{0.15cm}
    {\color{courseNavy}\rule{\linewidth}{1.4pt}\par}
    \vspace{0.35cm}
    {\Large\bfseries INDIVIDUAL TERM PROJECT\par}
    {\large\bfseries SUBMISSION FORM\par}
  \end{center}

  \vspace{0.45cm}
  \begin{projectbox}{courseNavy}
    \textbf{PROJECT TITLE}\par
    \vspace{2pt}
    {\large\bfseries PURPLE-LOOP --- A Dynamic Multi-Agent Orchestration Architecture for
    Adversarial Security Evaluation (Red vs Blue vs Audit)}
  \end{projectbox}

  \small
  \begin{tabularx}{\linewidth}{@{}L{0.30\linewidth}Y@{}}
    \textbf{Full Name} & Mohammed Mustafa Kareem \\
    \textbf{Student ID} & 230208878 \\
    \textbf{Department / Program} & Software Engineering \\
    \textbf{Email Address} & \href{mailto:230208878@ostimteknik.edu.tr}{230208878@ostimteknik.edu.tr} \\
    \textbf{GitHub / Project URL} & {\footnotesize\href{https://www.myskillos.com/project/490fe3ca-64b6-4831-b7e5-17a67adf7724}{myskillos.com/project/490fe3ca-64b6-4831-b7e5-17a67adf7724}} \\
    \textbf{Demo / Video URL (if available)} & Not provided \\
    \textbf{Submission Date} & 19 August 2026 \\
  \end{tabularx}

  \vfill
  \begin{center}
    \begin{tcolorbox}[width=0.88\linewidth,colback=courseBlue!6,colframe=courseBlue,
      arc=2mm,left=2mm,right=2mm,top=1mm,bottom=1mm]
      \centering
      {\normalsize\bfseries PROJECT SUBMISSION URL}\par
      {\Large\bfseries\href{https://myskillos.com/proje}{myskillos.com/proje}}\par
      \footnotesize Upload this form and the project files to the specified address.
    \end{tcolorbox}
    {\footnotesize 2026 Summer School -- Dr. Ufuk ASIL \& Dr. Maksat ATAGOZIEV}
  \end{center}
\end{titlepage}

\pagenumbering{arabic}

% ================================================================
\section*{Brief Description}
\addcontentsline{toc}{section}{Brief Description}

\begin{projectbox}{courseBlue}
PURPLE-LOOP is a dynamic multi-agent orchestration in which an attacker agent, a defender
agent and an independent auditor argue over a shared state object until the auditor
certifies convergence or the round budget is exhausted. Control flow is decided at runtime
by two LLM judgments --- the auditor's verdict and the router's routing decision. The
deterministic code only \emph{validates} moves; it never chooses them.
\end{projectbox}

\tableofcontents
\newpage

% ================================================================
\section{Project Summary}
\sectionintro{Provide a general overview of the project in a few short paragraphs.}

\qa{The problem the project aims to solve}{%
A classic security review is static: it walks a checklist, matches patterns with
\texttt{if/else}, and emits a score. Such a pipeline cannot reason about novel attack
composition, cannot argue with itself, and cannot decide that a particular question
deserves a second look. Its output is also unfalsifiable --- a defence is recorded as
effective because it was proposed, never because it was tested.

PURPLE-LOOP replaces the checklist with a structured disagreement. The core problem it
solves is making the \emph{control flow itself} a reasoning product: which agent speaks
next, and what it is asked to do, is derived from the current epistemic state of the
engagement rather than from a fixed script.}

\qa{The purpose of the project and its target users}{%
The purpose is to evaluate a described system architecture through adversarial debate and
to produce a security posture assessment whose every claim is traceable to a recorded
verdict. Target users are software engineering and security students running architecture
reviews as a tabletop exercise, and engineers who want a structured second opinion on a
design document before implementation. It is a teaching and design-review instrument, not
a penetration-testing tool.}

\qa{Brief project summary: How does the system work overall, and what does it produce?}{%
The user supplies a free-text description of a (fictional) system. A \textbf{Router}
(chief) agent inspects a compact digest of shared state and selects the next agent from a
formally computed set of legal moves, writing that agent a specific mandate.
\textbf{Recon} decomposes the target into assets and trust boundaries. \textbf{Red Team}
produces attack hypotheses in two lanes --- software/logic and social engineering ---
which pass through a \textbf{Safety Guardian} filter before entering state. \textbf{Blue
Team} answers with layered countermeasures. The \textbf{Security Auditor} issues one of
six verdicts per finding/mitigation pair; a \texttt{CHALLENGE} verdict sends the
defender's own patch back to the attacker as the next target. The loop ends on
convergence or budget exhaustion, after which the auditor authors the final report.

Output: a Markdown engagement report containing an agent-authored executive summary, a
0--100 security score with weighted breakdown, the full routing history, the disagreement
history, a per-agent calibration ledger, the safety-guardian log, and a complete
\texttt{final\_state.json} that allows any turn to be replayed.}

\begin{tabularx}{\linewidth}{@{}L{0.34\linewidth}Y@{}}
\textbf{Application domain} & Cybersecurity --- threat modelling and architecture review (tabletop, fictional targets) \\
\textbf{Languages / technologies used} & Python 3.10+ (standard library only), Markdown/YAML agent specifications, JSON Schema, LaTeX \\
\textbf{Models or services used} & Claude (\texttt{claude-sonnet-5}) through the Anthropic Messages API in \texttt{--live} mode; a deterministic mock brain for offline, key-free execution \\
\textbf{Current project status} & Complete and runnable. 10/10 architectural self-tests pass; a full 12-round engagement executes end to end \\
\end{tabularx}

% ================================================================
\section{Input, Output, and Usage}
\sectionintro{Summarize what the system receives, what it produces, and how it is used.}

\qa{Primary inputs}{%
(1) \texttt{target\_spec} --- a free-text description of a fictional system architecture:
components, data flows, stated controls and known constraints. The bundled example is
\texttt{sample\_target.md}, a fictional fintech named ``NovaPay''.
(2) \texttt{max\_rounds} --- the turn budget (6--12 recommended). No agent may extend it.}

\qa{Primary outputs}{%
(1) \texttt{sample\_run\_report.md} --- the engagement report: agent-authored executive
summary, weighted score breakdown, routing history with the legal moves available at each
turn, findings table, disagreement history, calibration ledger, guardian log, outstanding
risk.
(2) \texttt{final\_state.json} --- the complete \texttt{EngagementState} object, including
the append-only patch history for replay.
(3) A single composite security score on 0--100 with a letter grade.}

\qa{Brief usage scenario: Main steps from the user request to the result}{%
\begin{enumerate}[nosep]
  \item The user writes or selects an architecture description and sets a round budget.
  \item The orchestrator initialises \texttt{EngagementState} and calls the Router.
  \item The Router returns \texttt{next\_agent} plus a mandate; the orchestrator verifies
        the choice lies inside the legal-move set and runs that agent against its state
        projection.
  \item Attacker output is routed through the Safety Guardian; the resulting patch is
        validated against five invariants and applied or rejected with a logged reason.
  \item The scorecard and convergence metrics are recomputed; the loop repeats from step 3.
  \item On convergence or budget exhaustion the Auditor authors the final report, which is
        written to disk together with the final state.
\end{enumerate}}

\qa{A sample input and its corresponding output}{%
\textbf{Input (excerpt from \texttt{sample\_target.md}):} ``Support desk --- outsourced to a
BPO partner. Agents verify callers using name, date of birth, and the last four digits of
the registered phone number, then may trigger a password reset from the admin console.''

\textbf{Output (excerpt from the engagement report):} Finding \texttt{F1} (lane: social) ---
\emph{``Support desk identity verification relies on knowledge that is not secret.''}
Likelihood 4 / Impact 4, confidence 0.67. Round 3: Blue proposes mitigation \texttt{M1}
with a preventive and a containment layer only. Round 4: the Auditor returns
\texttt{INSUFFICIENT} --- \emph{``No detection layer: if prevention fails the event is
invisible.''} Round 6: Blue submits \texttt{M4} (version 2) with a detection layer added.
Round 7: verdict \texttt{SOUND}. Final engagement score: \textbf{82.6/100, grade B}.}

% ================================================================
\section{System Design}
\sectionintro{Describe the architecture only to the extent that it applies to your project.}

\qa{Why did you choose this approach? Briefly justify the choice of an agent, workflow, or traditional software}{%
A fixed workflow (Recon $\rightarrow$ Red $\rightarrow$ Blue $\rightarrow$ Audit $\rightarrow$ end)
was rejected because it wastes turns on solved areas and starves contested ones, and
because it can never revisit a defence once proposed. Traditional software was rejected
because the judgments involved --- ``is this attack reachable given these trust
boundaries?'', ``does this patch actually close the finding?'' --- are not expressible as
rules over a finite feature set.

The agent approach was chosen specifically so that \emph{routing} could be a reasoning
step. Two LLM decisions drive control flow: the Auditor's verdict on each
finding/mitigation pair, and the Router's choice of the next agent and its mandate.
Determinism is retained where it belongs: the legal-move set and the five state
invariants are computed in code, so the LLM exercises judgment inside a formally bounded
space rather than an unbounded one.}

\qa{The system's main components and their responsibilities}{%
Six agents plus a deterministic orchestration harness. The Router (chief) selects the next
agent and writes its mandate. Recon builds the asset and trust-boundary inventory. Red
Team generates attack hypotheses in two lanes. Blue Team produces three-layer
countermeasures. The Security Auditor scores, contests, routes and reports. The Safety
Guardian filters every attacker output. The harness owns state: it validates patches
against the invariants, computes the legal-move set, recomputes the scorecard, and grades
each routing prediction. Red and Blue never communicate directly --- all inter-agent
communication is mediated by shared state, which makes every turn inspectable and
replayable.}

\qa{The main data or task flow between components}{%
A single JSON object, \texttt{EngagementState}, is the only channel. Each agent receives a
\emph{projection} of it (only entitled fields, with prior-round findings compacted to a
headline) and returns a \emph{patch}, never a full rewrite. Patches are applied by the
orchestrator under five invariants: findings are append-only; every mitigation must
reference a live finding; the auditor may not author findings or mitigations; only the
guardian may write attacker content; no agent may extend the round budget. Every applied
or rejected patch is recorded in \texttt{patch\_history}.}

\begin{projectbox}{courseTeal}
\textbf{Architecture diagram:}\par
\vspace{2mm}
\IfFileExists{architecture.png}{\includegraphics[width=\linewidth]{architecture.png}}{}
\end{projectbox}

% ================================================================
\section{Agents, Tools, and Integrations}
\sectionintro{If your project does not use agents or external tools, you may write ``not used'' in the relevant rows.}

\begin{tabularx}{\linewidth}{@{}L{2.9cm}Y Y@{}}
\toprule
\textbf{Component / agent / tool} & \textbf{Brief responsibility} & \textbf{Information received or produced} \\
\midrule
\texttt{orchestrator-router} (Chief) &
Chooses the next agent and writes its mandate; may choose only from the legal-move set &
Receives a compact state digest and the legal moves; produces \texttt{next\_agent}, \texttt{mandate}, \texttt{lane}, rationale and a falsifiable prediction \\

\texttt{recon-analyst} &
Decomposes the target; surfaces the human and process assets the specification omitted &
Receives \texttt{target\_spec}; produces \texttt{asset\_inventory}, \texttt{trust\_boundaries}, \texttt{coverage\_map}, \texttt{unknowns} \\

\texttt{red-team} &
Attacker. Generates abstract attack hypotheses in the logic and social lanes; attacks deployed patches on a \texttt{CHALLENGE} mandate &
Receives assets, boundaries, existing mitigations and verdicts; produces \texttt{findings[]} with ATT\&CK IDs, likelihood, impact, confidence, novelty \\

\texttt{blue-team} &
Defender. Produces preventive, detective and containment layers per finding and declares the assumptions each control rests on &
Receives findings and verdicts; produces \texttt{mitigations[]} with \texttt{assumptions}, \texttt{does\_not\_stop}, residual-risk claim \\

\texttt{security-auditor} &
Independent judge; scores, contests, instructs, decides convergence, authors the report &
Receives the full state except \texttt{target\_spec}; produces \texttt{audit\_verdicts[]}, \texttt{scorecard}, \texttt{calibration\_ledger}, \texttt{final\_report} \\

\texttt{safety-guardian} &
Filters every attacker output before the state write: PASS, REDACT or REJECT &
Receives candidate findings; produces sanitized findings and \texttt{guardian\_log} \\
\bottomrule
\end{tabularx}

\qa{APIs, data sources, MCP, Skills, models, or other integrations used (if any)}{%
Anthropic Messages API (\texttt{claude-sonnet-5}) in \texttt{--live} mode; no key is
required for the default mock mode. The agent team is packaged as a Skill in the
Myskillos / Claude layout (\texttt{CLAUDE.md} plus \texttt{.claude/agents/*.md}), generated
from the prompt document by \texttt{python purple\_loop.py --build-package} so the report
and the uploaded skill cannot drift apart. Conceptual references only, with no runtime
dependency: the MITRE ATT\&CK technique vocabulary and the STRIDE threat categories. No
external data sources, databases or network calls are used at run time.}

\qa{Important decisions, permission boundaries, or operations requiring human approval (if any)}{%
Three boundaries are enforced structurally rather than by convention. (1) The Safety
Guardian sits between the attacker and shared state; attacker output that has not been
cleared is rejected by the patch validator, so the filter cannot be bypassed even by a
malformed agent response. (2) Separation of powers: the auditor is structurally prevented
from authoring findings or mitigations, so the judge cannot manufacture the evidence it
scores. (3) Human escalation: when no legal move would meaningfully advance the engagement
--- for example when the target description is too thin to support further analysis ---
the Router returns \texttt{next\_agent: "HUMAN"} with one specific question and the
engagement pauses rather than fabricating depth. The round budget is likewise a human-only
parameter; no agent can extend it.}

% ================================================================
\section{Implementation and Execution}

\begin{tabularx}{\linewidth}{@{}L{0.35\linewidth}Y@{}}
\textbf{Source code URL} & Included in this submission (\texttt{implementation/}). Myskillos project: \href{https://www.myskillos.com/project/490fe3ca-64b6-4831-b7e5-17a67adf7724}{myskillos.com/project/490fe3ca…} (agent team exported to \texttt{myskillos\_export/}) \\
\textbf{Main execution command} & \texttt{python purple\_loop.py --target sample\_target.md} \\
\textbf{Required core dependencies} & Python 3.10+ with the standard library only. Optional: \texttt{anthropic} (needed only for \texttt{--live}) \\
\textbf{Important files or folders} & \pl{} (orchestrator), \texttt{02\_AGENT\_PROMPTS.md} (all six system prompts), \texttt{03\_STATE\_SCHEMA.json}, \texttt{04\_SCORING\_RUBRIC.md}, \texttt{myskillos\_package/} \\
\textbf{Execution environment} & Any OS with Python 3.10+; no API key, network access or GPU required for the default run \\
\end{tabularx}

\qa{Brief installation and usage steps for running the project}{%
\begin{enumerate}[nosep]
  \item Extract the submission archive and open the \texttt{implementation/} folder.
  \item \texttt{python purple\_loop.py --self-test} --- asserts the ten architectural
        invariants. Expected: \texttt{10/10 checks passed}.
  \item \texttt{python purple\_loop.py --target sample\_target.md} --- runs a full
        engagement and writes \texttt{sample\_run\_report.md} and
        \texttt{final\_state.json}. No API key needed.
  \item Optional: \texttt{export ANTHROPIC\_API\_KEY=...} then add \texttt{--live} to drive
        the same orchestration with real LLM agents.
  \item Optional: \texttt{python purple\_loop.py --build-package} regenerates the
        Myskillos-importable agent team from the prompt document.
\end{enumerate}}

\qa{A significant technical problem encountered and the solution applied}{%
\textbf{Problem --- infinite adversarial escalation.} The \texttt{CHALLENGE} verdict is the
mechanism that makes the system a loop rather than a pipeline: the defender's patch becomes
the attacker's next target. In the first working version this never terminated. Each
successful bypass was recorded as a \emph{new} finding, which the defender then mitigated,
which the auditor then challenged, which the attacker then bypassed. The engagement
consumed its entire budget on a single ancestral finding and the score oscillated instead
of converging.

\textbf{Solution --- make ``no bypass found'' a first-class, scored outcome.} Two changes.
First, a reworked (version 2) layered control resists the bypass attempt, and the attacker
is required to report that failure honestly through its \texttt{self\_rejected} channel
rather than manufacturing a weaker finding to appear productive. Second, the auditor
re-audits the challenged pair once the challenge has been answered: if a bypass was
produced the verdict is \texttt{INSUFFICIENT}, and if none was produced the verdict is
\texttt{SOUND} with a \textbf{durability bonus} of $+3$. The bonus is the only way in the
rubric to earn points for a defence being \emph{tested} rather than merely asserted, so
the termination fix and the scoring model reinforce each other instead of competing. Turn
8 of the bundled run is exactly this case: the attacker produces nothing and the defender
scores for it.}

% ================================================================
\section{Tests and Results}
\sectionintro{Describe the basic checks performed to demonstrate that the project works.}

\begin{tabularx}{\linewidth}{@{}C{0.7cm}Y Y L{2.0cm}@{}}
\toprule
& \textbf{Trial / test} & \textbf{Result obtained} & \textbf{Status} \\
\midrule
1 & Invariant suite: \texttt{--self-test} asserts all five state invariants, the verdict vocabulary, guardian redaction, prompt loading and state projection & \texttt{10/10 checks passed}; attempts to bypass the guardian or to let the auditor author findings raise \texttt{InvariantError} & Pass \\
2 & Full engagement: 12-round run against \texttt{sample\_target.md} in mock mode & Ran to completion; 3 findings, 6 mitigations, 8 verdicts, 0 rejected patches; final score \textbf{82.6/100 (B)} & Pass \\
3 & Non-pipeline assertion: executed agent sequence compared against the fixed cycle Recon/Red/Blue/Audit & Sequence diverges from the fixed cycle; Red is called three times with three different mandates & Pass \\
4 & Safety Guardian: a candidate finding containing operational text is submitted to the filter & Action \texttt{REDACT}; the operational phrasing is replaced by the vulnerability class and the ATT\&CK reference & Pass \\
5 & Myskillos platform dry run: the six-agent team was built on the visual editor, saved and tested & \emph{``No issues found --- the structure is consistent''}; export reports full compatibility, 8 files & Pass \\
\bottomrule
\end{tabularx}

\qa{Overall result: To what extent did the project achieve its objective?}{%
The objective --- an orchestration whose control flow is a reasoning product rather than a
script --- was achieved and is machine-verifiable rather than merely claimed. Three moments
in the bundled run are unreachable in a static pipeline: round 5, which exists only because
a verdict sent control \emph{backwards} from judge to attacker; round 8, where the attacker
correctly produces nothing and the defender earns a durability bonus for surviving; and
rounds 9--11, where a defender that regressed was caught twice by the same rubric. All
three technical expectations of the course are implemented and covered by the self-test
suite. The scoring model, the six-verdict vocabulary and the calibration ledger together
make each score traceable to a specific recorded judgment.}

\qa{Known limitations or incomplete aspects}{%
\begin{itemize}[nosep]
  \item \textbf{Routing bias (measured, not hidden).} The policy prioritises answering
        outstanding verdicts over exploring untouched coverage, so a run with many
        challenge loops can exhaust its budget having probed only one lane. The bundled
        run ends at 0.30 attack-surface coverage. The fix is to reserve a fraction of the
        budget for exploration; it was left visible because a routing policy with a
        demonstrated failure mode is more informative than one never measured.
  \item \textbf{Mock brain.} The offline stand-in exercises the control flow faithfully,
        but its security \emph{content} is fixed text; only \texttt{--live} produces
        genuine agent reasoning.
  \item \textbf{Not empirical.} LLM agents are not penetration testers. Scores are
        relative evaluations of an architecture \emph{description}, never measurements of
        a running system.
  \item \textbf{Heuristic thresholds.} The 0.25 novelty floor and 3.0 score-delta epsilon
        were chosen by judgment and would need tuning against real engagements.
  \item Adversarial self-critique reduces but does not eliminate agreement bias; the
        novelty floor and lane switching are mitigations, not proofs.
\end{itemize}}

% ================================================================
\section{Security and Responsible Use}
\sectionintro{Evaluate only the items relevant to your project.}

\begin{tabularx}{\linewidth}{@{}C{1.0cm}Y@{}}
\toprule
\textbf{OK} & \textbf{Check} \\
\midrule
\yes & Secret keys and passwords were not included in the source code. \\
\yes & User data was processed only to the extent necessary. \\
\yes & Risky or irreversible operations were restricted. \\
\yes & Sources and third-party content used were identified. \\
\bottomrule
\end{tabularx}

\qa{Project-specific security, ethics, privacy, or copyright concerns and the measures taken}{%
The central concern in a Red-versus-Blue simulation is that an attacker agent could be
induced to emit operationally useful attack content. The measure taken is architectural
rather than advisory: a dedicated \texttt{safety-guardian} agent inspects every attacker
output before it can enter shared state, and the patch validator rejects any attacker
content that has not been cleared, so the filter cannot be skipped by a malformed
response. Attack reasoning is therefore constrained to vulnerability \emph{classes}, MITRE
ATT\&CK technique identifiers and kill-chain narratives; exploit code, payloads,
command lines, malware, evasion tradecraft and ready-to-send phishing text are out of
scope by construction. The guardian is instructed to prefer redaction over rejection so
analytical value is preserved, but the code and malware rules are reject-only.

The evaluation target is fictional (``NovaPay'') and the guardian additionally replaces any
real, named organisation, product version or person with a generic placeholder, so the
system cannot be pointed at a real organisation through the input alone. This constraint
deliberately reduces attack specificity: the project accepts a small loss of analytical
fidelity for a safety property treated as non-negotiable.

No personal data is processed and no credentials appear in the source. The API key, when
\texttt{--live} is used, is read from the environment and never written to state, logs or
the report. Third-party content is limited to the conceptual ATT\&CK and STRIDE
vocabularies, which are cited in the references and not redistributed.}

% ================================================================
\section{Conclusion and Future Work}

\qa{The most important results and learning outcomes obtained from the project}{%
The main result is that adversarial multi-agent orchestration can be made
\emph{auditable}. Because every routing decision carries a rationale and a falsifiable
prediction, and every verdict carries a rationale that cites state, the final score is
defensible line by line rather than being an opaque number.

The main learning outcome was that the difficult part of multi-agent design is not
prompting the agents --- it is deciding what the deterministic layer must own. Placing the
legal-move computation, the state invariants and the scoring arithmetic in code, and
leaving only judgment to the model, is what made the system both non-trivial and
debuggable. A second lesson came from the termination bug: self-critiquing agent teams
fail in two opposite directions --- agreeing too early, and escalating forever --- and both
needed an explicit countermeasure, the novelty floor for the first and the durability
bonus for the second.}

\qa{Improvements that could be made with more time}{%
\begin{itemize}[nosep]
  \item \textbf{Budgeted exploration.} Reserve a fraction of the round budget for coverage
        so challenge loops cannot starve an entire attack lane --- the one measured
        weakness of the current routing policy.
  \item \textbf{Parallel Red lanes.} Run the logic and social lanes concurrently and let
        the Router arbitrate between their findings, rather than switching sequentially.
  \item \textbf{Calibration over many runs.} The confidence multipliers currently reset
        each engagement; persisting them across runs would let agent reliability be
        learned rather than assumed.
  \item \textbf{Threshold tuning.} Fit the novelty floor and convergence epsilon against a
        corpus of real architecture reviews instead of choosing them by judgment.
  \item \textbf{Human-in-the-loop review.} Expose the \texttt{HUMAN} escalation as an
        interactive prompt so an operator can supply missing architecture detail mid-run.
\end{itemize}}

\qa{Core course concepts applied in the project}{%
\textbf{State management} --- a single JSON blackboard with patch-based writes, an
append-only history, per-agent read projections, context compaction and five enforced
invariants. \textbf{Reflection and feedback loops} --- three nested loops: intra-agent
draft/critique/revise, a six-verdict adversarial cross-critique in which the
\texttt{CHALLENGE} verdict re-routes control backwards, and engagement-level convergence
with cross-round confidence calibration. \textbf{Dynamic routing} --- an LLM router
operating over a formally constrained legal-move set, generating mandates rather than mere
destinations, and making falsifiable predictions that are graded the following round.
Alongside these: the orchestrator--worker pattern, separation of powers between generator
and judge, guardrails as an architectural component, and structured JSON contracts between
agents.}

% ================================================================
\section{References and Submission Checklist}

\qa{Documents, articles, datasets, models, open-source projects, or other resources used}{%
\begin{itemize}[nosep]
  \item MITRE ATT\&CK Enterprise technique taxonomy --- \url{https://attack.mitre.org}
        (used as a vocabulary for referring to techniques abstractly; not redistributed).
  \item STRIDE threat categorisation (Microsoft) --- used as the structure of the
        \texttt{coverage\_map}.
  \item Anthropic Claude documentation and Messages API --- \url{https://docs.claude.com}
        (model \texttt{claude-sonnet-5}, used in \texttt{--live} mode).
  \item Myskillos agent platform --- \url{https://myskillos.com} (skill packaging format:
        \texttt{CLAUDE.md} plus \texttt{.claude/agents/*.md}).
  \item SENG 456 course materials on agent orchestration, state management and dynamic
        routing.
  \item AI assistance: Claude was used to help design and implement the orchestration, as
        explicitly permitted by the project brief. All architectural decisions, the
        verdict vocabulary, the scoring model and the final content of this report were
        reviewed and are owned by the author.
\end{itemize}}

\begin{tabularx}{\linewidth}{@{}C{1.0cm}Y@{}}
\toprule
\textbf{OK} & \textbf{Submission item} \\
\midrule
\yes & A completed PDF version of this form is included. \\
\yes & The source code or project files are accessible. \\
\yes & Installation and execution instructions are included. \\
\yes & A test result or sample usage result is included. \\
\yes & No secret keys, passwords, or unauthorized personal data are included. \\
\yes & Demo, video, or supplementary file links work, if provided. \\
\bottomrule
\end{tabularx}

% ================================================================
\section{Student Declaration}

\begin{projectbox}{courseRed}
I declare that I prepared this project individually; identified the sources,
third-party code, and artificial intelligence assistance I used; and ensured that
the submission files contain no secret keys or unauthorized personal data.
\end{projectbox}

\vspace{6mm}
\begin{tabularx}{\linewidth}{@{}L{0.25\linewidth}Y@{}}
\textbf{Full Name} & Mohammed Mustafa Kareem \\
\textbf{Student ID} & 230208878 \\
\textbf{Date} & 19 August 2026 \\
\textbf{Signature} & \rule{\linewidth}{0.4pt} \\
\end{tabularx}

\end{document}
